% Generated draft: requires mathematical and source review.
\documentclass[11pt]{article}
\usepackage[margin=25mm]{geometry}
\usepackage{fontspec}
\setmainfont{Latin Modern Roman}
\usepackage{amsmath,amssymb,mathtools,unicode-math}
\setmathfont{Latin Modern Math}
\usepackage{enumitem,needspace,xurl,hyperref}
\setlength{\parindent}{0pt}
\setlength{\parskip}{5pt}
\setlength{\emergencystretch}{3em}
\newcommand{\N}{\mathbb{N}}
\newcommand{\Z}{\mathbb{Z}}
\newcommand{\Q}{\mathbb{Q}}
\newcommand{\R}{\mathbb{R}}
\newcommand{\C}{\mathbb{C}}
\newcommand{\F}{\mathbb{F}}
\newcommand{\A}{\mathbb{A}}
\newcommand{\PP}{\mathbb{P}}
\newcommand{\E}{\mathbb{E}}
\newcommand{\eps}{\varepsilon}
\newcommand{\dd}{\,\mathrm{d}}
\begin{document}
\section*{UnsolvedMath ID: 800700}
\textbf{Source problem code:} catalogue entry
\subsection{Definitions and mathematical statement}
Let \(F\) be a field. A central simple \(F\)-algebra is a finite-dimensional associative unital \(F\)-algebra that is simple and has center \(F\). Its degree is \(\deg(A)=\sqrt{\dim_F A}\). The Brauer group \(\operatorname{Br}(F)\) consists of equivalence classes of central simple \(F\)-algebras, where \(A\) and \(B\) are equivalent if \(M_r(A)\cong M_s(B)\) as \(F\)-algebras for some positive integers \(r,s\). The group operation is induced by tensor product over \(F\).

For a positive integer \(n\), define the subgroup
\[
H_n(F)=\bigl\langle [A]\in\operatorname{Br}(F): A\text{ is a central simple }F\text{-algebra with }\deg(A)=n\bigr\rangle.
\]
The supplied source asks, as a weaker question,
\[
\operatorname{Br}_p[n](F)\stackrel{?}{=}H_n(F).
\]
This equality transcribes the question under the usual interpretation of generation inside the Brauer group. The supplied record does not define \(\operatorname{Br}_p[n](F)\), specify the meaning or range of \(p\), give the permitted values of \(n\), or state any restrictions on \(F\). The field and positive-integer conventions above define the right-hand side; they do not resolve the missing source hypotheses. Thus the display is an unresolved formulation template, not a fully quantified conjecture. The stronger question to which the source compares it is also absent.
\subsection{Short English statement}
The source asks whether the group denoted by \(\operatorname{Br}_p[n](F)\) is generated by Brauer classes of central simple \(F\)-algebras of degree \(n\). Its notation and parameter restrictions are missing from the supplied excerpt, so a complete quantified statement requires further source context.
\subsection{Sources}
\begin{itemize}
\item Local UnsolvedMath registry: supplied complete source record with permanent ID \(800700\), exact source problem code \texttt{catalogue entry}, titled \emph{Open Problems on Central Simple Algebras --- Problem 10}. The registry attributes the problem to Auel, Brussel, Garibaldi, and Vishne, with proposed year \(2010\). Its \texttt{open} status is supplied registry metadata, not an independently verified mathematical or literature assessment.
\item Auel, Brussel, Garibaldi, and Vishne, \emph{Open Problems on Central Simple Algebras} (2010), Problem 10. Supplied source URL: \url{https://arxiv.org/abs/1006.3304}. The background records access on \(2026\text{-}10\text{-}16\) and extraction method \texttt{arxiv\_latex}; these are supplied provenance details, not a claim of independent source consultation.
\end{itemize}
\end{document}
